\section{引言}
\label{sec:intro}

记忆已成为现代大语言模型（LLM）智能体的核心组件，使其能够超越单个上下文窗口保留并复用信息，从而支持长程交互、个性化与知识密集型推理~\citep{hu2025memory}。
然而，既有工作大多聚焦于离线、查询无关的记忆构建：它们以固定方式预处理、压缩或索引历史上下文，而不考虑下游查询~\citep{kang2025memory, fang2025lightmem, zhong2024memorybank, xu2025mem, packer2023memgpt, chhikara2025mem0}。
这种“\textit{一次构建，始终使用}”的范式既可能浪费资源，也较为脆弱：无论具体查询需要什么，它都会消耗计算，同时还可能遗漏某些查询所需的关键信息。

一种直观的替代方案是按需记忆抽取，即根据当前查询在运行时触发计算。
这种灵活性也有代价：记忆处理被推迟到运行时，使成本与延迟成为首要问题。
在实际应用中，工业级 LLM 系统越来越多地提供显式且通常\textit{分档}的计算控制，例如“思考”模式~\citep{singh2025openai}、推理等级~\citep{openai-codex}或更重型的模型选项~\citep{claude-code}；这反映出系统需要在质量与运行时成本之间取得平衡。
由此，智能体记忆面临一个关键问题：\textit{如何为运行时记忆抽取提供显式且可控的性能--成本权衡？}

遗憾的是，在运行时智能体记忆中实现性能--成本权衡存在根本性困难。
现有权衡机制大多在离线阶段工作，而按需记忆会把这些决策推迟到运行时，使每个查询都需要在质量与成本之间作出选择。
这暴露出两个核心问题。
第一，\textit{预算应施加在哪里？}
现有系统常将记忆视为采用固定计算设置的单体流水线，导致权衡粒度粗糙且难以控制~\citep{kang2025memory, zhong2024memorybank}。
因此，需要为运行时场景定义合适的预算单元，即确定记忆抽取过程中的哪些模块应被分配预算，从而有针对性且有效地控制计算。
第二，\textit{预算应如何实现？}
既有工作很少为运行时记忆的权衡提供系统性指导，往往诉诸临时性的成本削减，或仅仅增加计算量~\citep{yan2025general}。
因此，即使预算方案已经确定，如何落实预算控制、哪些设计维度最能表达有意义的权衡，以及这些选择在不同预算区间下如何表现，仍不清楚。
要回答这些问题，就必须超越一次性的启发式方法或高计算量的简单升级，以更系统的视角研究运行时智能体记忆的性能--成本控制。

为解决上述挑战，我们提出 \textbf{BudgetMem}，一个运行时智能体记忆框架，通过按需记忆抽取实现显式、可控的性能--成本权衡。
BudgetMem 将运行时记忆抽取视为多阶段模块化流水线，并使计算能够在模块级得到控制。
具体而言，BudgetMem 通过统一的预算档位接口规范各模块的调用方式，使学习得到的路由器可以在保持整体抽取结构不变的同时，为每个模块选择不同预算档位。

在这一模块化骨架之上，BudgetMem 为每个模块提供三个预算档位，即 \textsc{低}/\textsc{中}/\textsc{高}，对应不同的质量--成本权衡。
我们通过三种互补的分层策略实现预算档位：\textit{实现分层}（改变模块实现）、\textit{推理分层}（改变推理行为）和\textit{容量分层}（改变模块的模型容量）。
为在这些分层选项中进行选择，BudgetMem 使用一个共享的轻量级路由器，在处理查询的过程中执行预算档位路由：路由器在每个模块处根据可用上下文，即查询与模块中间状态，选择一个档位。
我们采用带成本感知奖励的强化学习来训练路由器，在任务性能与记忆抽取成本之间进行权衡，从而形成可控的性能--成本行为。
这一设计既得到了一套实用、可控的运行时记忆系统，也支持在统一框架内比较不同的预算实现策略。

在 LoCoMo、LongMemEval 和 HotpotQA 上的实验表明，在性能优先的高预算设置下，BudgetMem 相比有竞争力的基线取得显著提升；随着预算收紧，它也呈现出清晰的性能--成本权衡曲线。
此外，我们的分析分离了不同预算分层策略的相对优势，为判断不同预算区间内哪种机制能带来最佳收益提供了依据。

本文贡献总结如下：
\begin{itemize}
\item 我们提出 \textbf{BudgetMem}，一种模块化运行时智能体记忆框架，通过具有预算档位的模块，对按需记忆抽取实施显式的性能--成本控制。
\item 我们提出预算档位路由，使用强化学习训练一个共享的轻量级路由器，在抽取过程中选择预算档位；并在统一框架内进一步实现三种互补的预算策略，即\textit{实现分层}、\textit{推理分层}和\textit{容量分层}。
\item LoCoMo、LongMemEval 和 HotpotQA 上的实验展示了强劲性能与清晰的性能--成本权衡；进一步分析还揭示了不同策略在各预算区间内何时能够取得最佳收益。
\end{itemize}
